% Generated draft: requires mathematical and source review.
\documentclass[11pt]{article}
\usepackage[margin=25mm]{geometry}
\usepackage{fontspec}
\setmainfont{Latin Modern Roman}
\usepackage{amsmath,amssymb,mathtools,unicode-math}
\setmathfont{Latin Modern Math}
\usepackage{enumitem,needspace,xurl,hyperref}
\setlength{\parindent}{0pt}
\setlength{\parskip}{5pt}
\setlength{\emergencystretch}{3em}
\newcommand{\N}{\mathbb{N}}
\newcommand{\Z}{\mathbb{Z}}
\newcommand{\Q}{\mathbb{Q}}
\newcommand{\R}{\mathbb{R}}
\newcommand{\C}{\mathbb{C}}
\newcommand{\F}{\mathbb{F}}
\newcommand{\A}{\mathbb{A}}
\newcommand{\PP}{\mathbb{P}}
\newcommand{\E}{\mathbb{E}}
\newcommand{\eps}{\varepsilon}
\newcommand{\dd}{\,\mathrm{d}}
\begin{document}
\section*{UnsolvedMath ID: 116}
\textbf{Source problem code:} catalogue entry
\subsection{Definitions and mathematical statement}
Let $p$ be an odd prime, let $n\geq1$ be an integer, and put $G=\F_p^n$. For a finite nonempty set $S$, write $\E_{s\in S}F(s)=|S|^{-1}\sum_{s\in S}F(s)$. Let $f:G\times G\to\C$ satisfy $|f(x,y)|\leq1$ for every $x,y\in G$.

Using the standard multiplicative-derivative convention, set
\[
 \Delta_{(h,h)}f(x,y)=f(x+h,y+h)\overline{f(x,y)}\qquad(h\in G).
\]
For $g:G\times G\to\C$, its box norm is defined by
\[
 \|g\|_{\square}^4=
 \E_{x,x',y,y'\in G}
 g(x,y)\overline{g(x',y)}\,\overline{g(x,y')}g(x',y').
\]
The right-hand side is real and nonnegative. Replacing the displayed derivative by its complex conjugate leaves its box norm unchanged, so the derivative orientation does not change the hypothesis. Suppose, for a positive threshold $\delta$ (necessarily at most $1$ when the hypothesis holds), that
\[
 \E_{h\in G}\|\Delta_{(h,h)}f\|_{\square}^4\geq\delta.
\]

\textbf{Source question.} Must $f$ correlate with a function of the displayed form
\[
 H(x,y)=a(x)b(y)c(x+y)(-1)^{q(x,y)}?
\]
Here $a,b,c$ are functions of the indicated single $G$-variables; their codomains and magnitude bounds are not stated in the record. If $H$ is a specified complex-valued function, normalized correlation would be measured by
\[
 \left|\E_{x,y\in G}f(x,y)\overline{H(x,y)}\right|.
\]
The source does not give the required positive lower bound, its dependence on $\delta$, or its uniformity in $n$ and $p$. In particular, writing a dimension-independent lower bound $\eta(\delta)>0$ would be an additional specification that must be checked against the intended theorem.

Most substantially, the source does not define the class or codomain of $q$, or the meaning of the phase $(-1)^{q(x,y)}$. For odd $p$, parity of a representative of an element of $\F_p$ is not an intrinsic field-valued operation. Therefore the displayed phase cannot silently be interpreted as an additive character or replaced by one; nor can $q$ silently be required to be quadratic. These unresolved choices prevent a complete precise inverse-theorem statement.

The supplied literature triage dated 2026-08-17 identifies this with Problem 55 of Ben Green's maintained collection, although its local code is catalogue entry. It reports that the inverse theorem is unknown and describes this as the simplest Gowers-type norm with an unknown inverse theorem. It also explicitly flags the same odd-characteristic phase ambiguity in the local import and the cited PDF. Those are dated source assessments, not an independently verified current status.
\subsection{Short English statement}
For a bounded complex function on $\F_p^n\times\F_p^n$, ask whether a large average box norm of its diagonal multiplicative derivatives forces correlation with a product depending separately on $x$, $y$, and $x+y$, multiplied by the displayed phase $(-1)^{q(x,y)}$. The source leaves $q$, that phase, and the quantitative correlation requirements undefined.
\subsection{Sources}
\begin{itemize}
\item UnsolvedMath local registry, data/unsolvedmath/problems.json, permanent ID 116, source code catalogue entry, ``Gowers Box Norms over Finite Fields.'' The local set is Ben Green's 100 Open Problems.
\item B. Green, \emph{100 Open Problems}, Problem 55, version described in the supplied review as retrieved 2026-08-17: \url{https://people.maths.ox.ac.uk/greenbj/papers/open-problems.pdf}. According to that review, the source gives the same formula, diagonal derivative and box configuration, and records the inverse theorem as unknown. No PDF was opened or independently checked in preparing this draft.
\item Ulam AI and the UnsolvedMath Contributors, \emph{UnsolvedMath: A Curated Collection of Open Mathematics Problems}; local record and its literature triage dated 2026-08-17. Dataset curation and original metadata are supplied under CC BY 4.0; referenced source material retains its own terms.
\end{itemize}
\end{document}
